\section{Peer review draft --- M2}\label{doc7-peer-review-draft-m2}

\textbf{Тэнгис · Nogoolin · 2026-09-21 · v0.2 · Бэлтгэсэн; илгээгээгүй}

\subsection{Reviewer-т зориулсан
тайлбар}\label{doc7-reviewer-ux442-ux437ux43eux440ux438ux443ux43bux441ux430ux43d-ux442ux430ux439ux43bux431ux430ux440}

Week 2-оос Document RAG Chatbot-оос Nogoolin руу шилжсэн. Week 1-ийн
ажлыг өөрчлөөгүй. Phase 0 баримтаас SRS-ийн шаардлага болон SDD-ийн
технологийн шийдвэрийг салгаж, таван шаардлагын draft ба traceability
бэлтгэлээ.

\textbf{Унших дараалал:}
\href{../../requirements/scope-charter.md}{Scope Charter} →
\href{wiki/persona-and-pivot.md}{persona нэмэлт} →
\href{../../requirements/srs-template.md}{SRS загвар} →
\href{../../requirements/requirements.md}{5 шаардлага} →
\href{../../requirements/traceability-matrix.md}{traceability} →
\href{../../architecture/sdd-template.md}{SDD}.

\subsection{Шалгах
асуултууд}\label{doc7-ux448ux430ux43bux433ux430ux445-ux430ux441ux443ux443ux43bux442ux443ux443ux434}

\begin{enumerate}
\def\labelenumi{\arabic{enumi}.}
\tightlist
\item
  Included/excluded/postponed нь ганцаараа хөгжүүлэх хүрээнд хангалттай
  тодорхой байна уу? Product 360° viewer ба нүүрний hero хоёр ялгагдаж
  байна уу?
\item
  SRS загвар IEEE 830-ын агуулгын бүтцийг хадгалж, ISO/IEC/IEEE
  29148-ийн requirement card, priority, source/owner, verification болон
  traceability-г зөв хослуулсан уу?
\item
  Mapping-ийн мөр бүр яг нэг төлөвтэй, n-a нь нэг өгүүлбэр үндэслэлтэй
  байна уу? \nolinkurl{complete} нь тухайн M2 артефакт бэлэн гэсэн
  утгатайг ойлгомжтой тайлбарласан уу?
\item
  FR-01/02, NFR-01\ldots03 бүрд оролт, үйлдэл, гаралт, pass/fail хязгаар
  хангалттай юу? Олон шалгуурыг M3-д дэд шаардлага болгох хэрэгтэй юу?
\item
  \textbf{W1 P1 → P-NG-01 → NFR-01}, \textbf{W1 P3 → P-NG-03 → NFR-03}
  гэсэн тайлбарласан холбоо нь төсөл сольсон нөхцөлд US-2.3-ын persona
  шаардлагыг хангах уу?
\item
  Ownership-ийг session-ээс авах дүрэм ойлгомжтой юу? Guest inquiry,
  customer history, admin inbox-ийг зөв ялгасан уу?
\item
  NFR-03-т all-skipped тестийг fail гэж үзэх шалгуур хангалттай юу?
  SRS/SDD-ийн хил зааг алдагдсан өгүүлбэр байна уу?
\end{enumerate}

\subsection{Санал бүртгэх
загвар}\label{doc7-ux441ux430ux43dux430ux43b-ux431ux4afux440ux442ux433ux44dux445-ux437ux430ux433ux432ux430ux440}

\begin{longtable}[]{@{}
  >{\raggedright\arraybackslash}p{(\linewidth - 6\tabcolsep) * \real{0.2000}}
  >{\raggedright\arraybackslash}p{(\linewidth - 6\tabcolsep) * \real{0.1900}}
  >{\raggedright\arraybackslash}p{(\linewidth - 6\tabcolsep) * \real{0.3500}}
  >{\raggedright\arraybackslash}p{(\linewidth - 6\tabcolsep) * \real{0.2600}}@{}}
\toprule\noalign{}
\begin{minipage}[b]{\linewidth}\raggedright
Reviewer / огноо
\end{minipage} & \begin{minipage}[b]{\linewidth}\raggedright
Requirement/section
\end{minipage} & \begin{minipage}[b]{\linewidth}\raggedright
Асуудал ба санал
\end{minipage} & \begin{minipage}[b]{\linewidth}\raggedright
Хариу / өөрчлөлт
\end{minipage} \\
\midrule\noalign{}
\endhead
\bottomrule\noalign{}
\endlastfoot
Pending & Pending & Санал хараахан ирээгүй & Pending \\
\end{longtable}

\subsection{Source interview ба
submission}\label{doc7-source-interview-ux431ux430-submission}

Source interview хийгдээгүй; owner нь Тэнгис, source нь repo-ийн файл ба
өөрийн ажиглалт. Хүлээн авагч болон Confluence/LMS холбоос байхгүй тул
2026-09-15-ны хэрэглэгчийн заавраар илгээлт/нийтлэлийг алгассан. Peer
review хийсэн, neighboring team-д илгээсэн гэж тэмдэглэхгүй. Дараа нь
хүлээн авах суваг тодорхой болбол Confluence Requirements хэсэгт repo
permalink холбож, submission evidence хадгалж болно.

\subsection{Богино
retrospective}\label{doc7-ux431ux43eux433ux438ux43dux43e-retrospective}

\textbf{Сайн болсон:} Phase 0-ийн их хэмжээний эхээс M2-д хэрэгтэй жижиг
хүрээ сонгож, кодтой тулгав. \textbf{Тодорхойгүй байсан:} SRS-ийн ``юу
хийх'' ба SDD-ийн ``яаж хийх'' өгүүлбэрүүд холилдсон; W1 persona нь өөр
төсөл байсан. \textbf{Дараагийн өөрчлөлт:} шаардлага бүрийг source →
observable result → verification дарааллаар бичиж, тестийн команд
амжилттай дууссан эсэхээс гадна тест үнэхээр ажилласан эсэхийг шалгана.

Chinchilla-ийн бүлэг 5-аас далд таамаг ба одоогийн дуусах нөхцөлийг
эхэлж тодруулах аргыг үргэлжлүүлэн хэрэглэнэ. Номын үйл явц ба энэ ажлын
хамгийн том зөрүү нь бодит хэрэглэгч/reviewer-ийн feedback хараахан
аваагүй байхад draft-ыг кодын шинжилгээнээс эхлүүлсэн явдал юм.
